\begin{center}
{\zihao{3}\heiti 通用神经网络处理器的结构感知切图与分层调度模型}\par\medskip
{\heiti 摘要}
\end{center}
多核神经网络处理器的执行效率受计算依赖、共享带宽与核内存储共同限制。针对三种逐层增加同步和缓存能力的运行场景，本文在保持原图语义与附件调度规则的条件下，建立以完工时间为主、兼顾额外搬运的切图与调度方法。

首先，将张量--操作二部图投影为操作依赖图，提取弱连通分量、层宽、双管道工作量与共享输入。对于无核间缓存同步的场景，采用完整分量装箱与结构切分相结合的有限构造，以任务工作量和依赖估计筛选方案，并对优先候选进行有限重排与合并。具备硬件同步后，将通信边界转为核心边界，以张量生命周期估计核内驻留风险，在同核复用与计算均衡之间选择。对于共享只读 FIFO Cache，进一步建立 DDR、L2 双带宽池事件模型，描述发射查询、完成插入与淘汰，并用核内换出估计修正候选时间。

冻结算法分别从原始输入独立生成方案，再由附件程序验收。100 图五核平均加速比分别为 3.8311、3.9522 和 4.1939。第二问的平均额外搬运较第一问减少 36.17\%；第三问固定同一方案开启 L2 后，平均耗时缩短 2.976\%，平均字节命中率为 26.45\%。三问已测正式方案均通过执行验收。独立新图检验同时发现第二问对输入表示和部分分层结构不够稳定，说明冷启动与有效的结构分析仍需配合独立性能验证。

本文通过有限结构构造和逐层细化的资源模型控制方案生成成本，用逐图对照区分通信减少、调度改变与缓存开关的效果。该方法在给定数据上兼顾求解速度与多核收益，后续改进重点为驻留时序近似和未见结构下的候选选择。

\textbf{关键词：}计算图切分；多核调度；张量生命周期；FIFO Cache；事件模型
